\section{Rapid relaxation induced by particle respawning at the system boundary}\label{sec:appendix}
In the stochastic simulations, particles absorbed at the outer boundary are respawned at the inner boundary in order to enforce flux conservation and maintain a stationary concentration profile.
Here, we provide a simple analytical argument demonstrating that this respawning mechanism not only produces the correct steady state, but also accelerates relaxation toward it.

To this end, we consider a one-dimensional diffusion problem on the interval $0 \le x \le L$,
\begin{equation}
    \frac{\partial c}{\partial t} = D \frac{\partial^{2} c}{\partial x^{2}}, \qquad c(0,t) = 0,
    \label{eq:1D_diff}
\end{equation}
and decompose the concentration into a steady-state profile and a transient deviation,
\begin{equation}
    c(x,t) = c_{\mathrm{st}}(x) + u(x,t).
    \label{eq:decomposition}
\end{equation}

The deviation $u(x,t)$ satisfies
\begin{equation}
    \frac{\partial u}{\partial t} = D \frac{\partial^{2} u}{\partial x^{2}}, \qquad u(0,t) = 0,
    \label{eq:u_eq}
\end{equation}
together with a boundary condition at $x = L$ that depends on how absorbed particles are handled.

\subsection{Case 1: Dirichlet boundary condition}\label{sec:case1}
If the concentration is fixed at the outer boundary, $c(L,t) = c_{0}$, then the deviation satisfies
\begin{equation}
    u(L,t) = 0.
    \label{eq:bc_dirichlet}
\end{equation}

Using separation of variables, $u(x,t) = X(x) T(t)$, the spatial eigenvalue problem becomes
\begin{equation}
    X'' + k^{2} X = 0, \qquad X(0) = X(L) = 0,
    \label{eq:eigen_dirichlet}
\end{equation}
with eigenfunctions $X_{n}(x) = \sin(k_{n} x)$ and eigenvalues
\begin{equation}
    k_{n} = \frac{n\pi}{L}, \qquad n = 1, 2, \ldots
    \label{eq:kn_dirichlet}
\end{equation}

The slowest decaying mode corresponds to $n = 1$, yielding a relaxation time
\begin{equation}
    \tau_{1} = \frac{1}{D(\pi/L)^{2}} = \frac{L^{2}}{\pi^{2} D}.
    \label{eq:tau1}
\end{equation}

\subsection{Case 2: Slope-matching boundary condition}\label{sec:case2}
In the respawning scheme used in the simulations, flux conservation is enforced by requiring that the diffusive flux entering the system at $x = 0$ matches the flux leaving at $x = L$.
This corresponds to a slope-matching condition,
\begin{equation}
    \left.\frac{\partial c}{\partial x}\right|_{x=0} = \left.\frac{\partial c}{\partial x}\right|_{x=L},
    \label{eq:slope_match}
\end{equation}
which implies for the deviation
\begin{equation}
    \left.\frac{\partial u}{\partial x}\right|_{x=0} = \left.\frac{\partial u}{\partial x}\right|_{x=L}.
    \label{eq:slope_match_u}
\end{equation}

The spatial eigenvalue problem now becomes
\begin{equation}
    X'' + k^{2} X = 0, \qquad X(0) = 0, \qquad X'(0) = X'(L),
    \label{eq:eigen_slope}
\end{equation}
with solutions $X_{n}(x) = \sin(k_{n} x)$ and eigenvalues determined by
\begin{equation}
    \cos(k_{n} L) = 1 \implies k_{n} = \frac{2\pi n}{L}, \qquad n = 1, 2, \ldots
    \label{eq:kn_slope}
\end{equation}

The slowest decaying mode corresponds to $n = 1$, yielding a relaxation time
\begin{equation}
    \tau_{2} = \frac{1}{D(2\pi/L)^{2}} = \frac{L^{2}}{4\pi^{2} D}.
    \label{eq:tau2}
\end{equation}

\subsection{Comparison and implication}\label{sec:comparison}
Comparing the two relaxation times,
\begin{equation}
    \tau_{2} = \frac{1}{4}\tau_{1},
    \label{eq:tau_ratio}
\end{equation}
we find that the slope-matching (respawning) boundary condition leads to relaxation that is four times faster than the standard Dirichlet boundary condition.
This acceleration arises because the lowest admissible spatial mode has twice the wavenumber, resulting in steeper gradients and faster diffusive decay.

This analysis demonstrates that respawning absorbed particles at the system boundary not only preserves the correct steady-state concentration profile, but also significantly enhances convergence toward that steady state.
As a result, the respawning scheme provides an efficient and well-controlled method for enforcing stationarity in stochastic diffusion simulations.
